% 04-vi-sao-gradient-bien-mat.tex ;  Why gradient vanishes in the copy task
%
% Scope: apply BPTT recurrence from ch01 to copy task; show how W_hh
% and tanh saturation combine to produce exponential decay.

\section{Vì sao gradient biến mất trong bài toán sao chép}
\label{sec:ch02-vi-sao}

Những con số ở mục trên đặt ra một câu hỏi: \emph{tại sao} gradient yếu đi 25
lần khi đi từ $t = 21$ về $t = 1$ trong
\tn{bài toán sao chép}{copy task} $T_{\text{mem}} =
10$? Câu trả lời nằm trong chính công thức BPTT mà chương~01 đã dẫn xuất.

\subsection{Sự lặp lại của $\mat{W}_{hh}$}

Nhìn lại phương trình~\eqref{eq:bptt-dh-recurrence}: gradient truyền từ
$\vect{h}_{t+1}$ về $\vect{h}_t$ qua hai con đường. Con đường thứ nhất; từ
output hiện tại qua $\mat{W}_{hy}\transpose$; chỉ đi một bước và không tích
lũy. Con đường thứ hai; từ tương lai qua $\mat{W}_{hh}\transpose$; \emph{lặp
lại} ở mọi bước lùi.

Hình~\ref{fig:ch02-rnn-gradient} minh họa điều này: gradient từ $L$ truyền
ngược qua từng hidden state, và mỗi lần nó đi qua một mũi tên $\vect{h}_{t+1}
\to \vect{h}_t$, nó bị nhân với $\mat{W}_{hh}\transpose$ thêm một lần nữa.

\begin{figure}[ht]
  \centering
  % ch02-rnn-unrolled-gradient.tex — RNN unrolled 5 steps with gradient arrows
%
% Shows the forward pass (left to right) and BPTT backward pass (right to left).
% Arrow thickness/opacity encodes gradient norm: thick+dark near output,
% thin+light far from output.

\begin{tikzpicture}[
    >=Stealth,
    node distance=1.8cm and 1.6cm,
    every node/.style={font=\footnotesize},
    % Forward arrows: standard
    fwd/.style={->, thick, black!70},
    % Gradient arrows: thickness varies by step
    grad strong/.style={->, line width=1.2pt, black!80},
    grad medium/.style={->, line width=0.7pt, black!50},
    grad weak/.style={->, line width=0.4pt, black!25},
    grad vweak/.style={->, line width=0.2pt, black!15},
  ]

  % Hidden states (top row)
  \node[draw, circle, minimum size=0.7cm] (h0) {$\vect{h}_0$};
  \node[draw, circle, minimum size=0.7cm, right=of h0] (h1) {$\vect{h}_1$};
  \node[draw, circle, minimum size=0.7cm, right=of h1] (h2) {$\vect{h}_2$};
  \node[draw, circle, minimum size=0.7cm, right=of h2] (h3) {$\vect{h}_3$};
  \node[draw, circle, minimum size=0.7cm, right=of h3] (h4) {$\vect{h}_4$};
  \node[draw, circle, minimum size=0.7cm, right=of h4] (h5) {$\vect{h}_5$};

  % Inputs (bottom left)
  \node[below=0.8cm of h1] (x1) {$\vect{x}_1$};
  \node[below=0.8cm of h2] (x2) {$\vect{x}_2$};
  \node[below=0.8cm of h3] (x3) {$\vect{x}_3$};
  \node[below=0.8cm of h4] (x4) {$\vect{x}_4$};
  \node[below=0.8cm of h5] (x5) {$\vect{x}_5$};

  % Outputs (top)
  \node[above=0.7cm of h1] (y1) {$\vect{y}_1$};
  \node[above=0.7cm of h2] (y2) {$\vect{y}_2$};
  \node[above=0.7cm of h3] (y3) {$\vect{y}_3$};
  \node[above=0.7cm of h4] (y4) {$\vect{y}_4$};
  \node[above=0.7cm of h5] (y5) {$\vect{y}_5$};

  % Forward arrows (input → hidden, hidden → hidden, hidden → output)
  \draw[fwd] (x1) -- (h1);
  \draw[fwd] (x2) -- (h2);
  \draw[fwd] (x3) -- (h3);
  \draw[fwd] (x4) -- (h4);
  \draw[fwd] (x5) -- (h5);

  \draw[fwd] (h0) -- (h1) node[midway, above] {$\mat{W}_{hh}$};
  \draw[fwd] (h1) -- (h2);
  \draw[fwd] (h2) -- (h3);
  \draw[fwd] (h3) -- (h4);
  \draw[fwd] (h4) -- (h5);

  \draw[fwd] (h1) -- (y1);
  \draw[fwd] (h2) -- (y2);
  \draw[fwd] (h3) -- (y3);
  \draw[fwd] (h4) -- (y4);
  \draw[fwd] (h5) -- (y5);

  % Loss
  \node[right=0.5cm of y5] (L) {$L$};
  \draw[fwd] (y5) -- (L);

  % Backward gradient arrows (below the hidden states, going right-to-left)
  % Gradient is strongest near the output, weakest far away
  \draw[grad strong, <-] (h4.south) to[bend left=30]
    node[below=0.15cm, font=\tiny] {$\partial L/\partial\vect{h}_5$} (h5.south);
  \draw[grad medium, <-] (h3.south) to[bend left=30]
    node[below=0.15cm, font=\tiny] {$\partial L/\partial\vect{h}_4$} (h4.south);
  \draw[grad weak, <-] (h2.south) to[bend left=30]
    node[below=0.15cm, font=\tiny] {$\partial L/\partial\vect{h}_3$} (h3.south);
  \draw[grad vweak, <-] (h1.south) to[bend left=30]
    node[below=0.15cm, font=\tiny] {$\partial L/\partial\vect{h}_2$} (h2.south);
  \draw[grad vweak, <-] (h0.south) to[bend left=40]
    node[below=0.2cm, font=\tiny] {$\approx \vect{0}$} (h1.south);

  % Label
  \node[below=1.5cm of h3, font=\small\itshape, black!60]
    {Mũi tên càng mảnh = gradient càng yếu};

\end{tikzpicture}

  \caption{RNN unrolled 5 bước với BPTT. Mũi tên đậm: gradient mạnh ở gần $L$.
    Mũi tên mảnh: gradient yếu dần khi lùi xa khỏi output. Độ dày mũi tên tỉ
    lệ với norm gradient tại bước đó.}
  \label{fig:ch02-rnn-gradient}
\end{figure}

Sau $T$ bước lùi, gradient tại $\vect{h}_1$ chứa tích của $T$ ma trận
$\mat{W}_{hh}\transpose$ (xen kẽ với các ma trận đường chéo từ đạo hàm
$\tanh$). Nếu $\|\mat{W}_{hh}\| < 1$, tích này tiến về 0 theo hàm mũ:

\[
\bigl\|\pd{L}{\vect{h}_1}\bigr\|
\approx \|\mat{W}_{hh}\|^T \cdot \bigl\|\pd{L}{\vect{h}_T}\bigr\|.
\]

\subsection{Vai trò của $\tanh$ saturation}

$\mat{W}_{hh}$ không phải là thủ phạm duy nhất. Đạo hàm của $\tanh$ tại
$\vect{h}_{t+1}$ là $1 - \vect{h}_{t+1}^2$. Khi một phần tử của hidden
state \tn{bão hòa}{saturation} (tiến gần đến $\pm 1$), đạo hàm của nó tiến gần
đến 0, và
gradient đi qua phần tử đó bị chặn lại.

Trong bài toán sao chép, RNN có xu hướng đẩy hidden state đến gần bão hòa
để ``ghi nhớ'' thông tin; các nghiên cứu sau này gọi đây là \enquote{storing
information at the saturation points of the activation function}. Nhưng chính
hành vi đó làm gradient yếu thêm: không chỉ $\mat{W}_{hh}\transpose$ làm nó
nhỏ đi, mà các hệ số $1 - h_{t,i}^2$ còn nhân thêm một hệ số $< 1$ nữa ở mỗi
bước.

\subsection{Không phải mọi RNN đều bị vanishing như nhau}

Kết quả từ \tn{bài toán cộng}{adding problem} cho thấy vanishing gradient không
phải là định mệnh:
với $d_h = 24$, gradient gần như không suy giảm ở $T = 50$. Điều này là vì
$\mat{W}_{hh}$ với khởi tạo $\mathcal{N}(0, 0.1)$ và $d_h = 24$ có trị riêng
lớn nhất đủ nhỏ, và vì bài toán cộng không ép mạng phải dùng đến bão hòa
$\tanh$.

Nhưng với $d_h = 12$ và $T = 100$, vanishing xuất hiện. Với $d_h = 32$ và
bài toán sao chép, nó xuất hiện ở $T_{\text{mem}} = 10$. Không có một ngưỡng
cố định; vanishing gradient là một hàm của ba biến: kích thước mạng, khởi
tạo, và áp lực từ bài toán.

Chương~3 sẽ cho bạn công cụ để tính toán chính xác khi nào vanishing xảy ra,
thông qua trị riêng lớn nhất của $\mat{W}_{hh}$ và
\tn{điều kiện đủ}{sufficient condition} của Pascanu và cộng sự.
